\section{Dualidad de lectura HMT: dualidad \texorpdfstring{\(T\)}{T}, bloque \texorpdfstring{\(243\)}{243}, código ternario de Golay y rutas de masa}
\label{sec:sucesor83-f083-dualidad-t}

La dualidad \(T\) de cuerdas se realiza como proyección continua de una
dualidad HMT anterior: el intercambio entre el valor visible y la memoria de
frontera. Todo entero admite la lectura
\[
n=9Q+r,
\]
donde \(r\) es el valor visible de la rueda nonádica y \(Q\) conserva la
memoria de cruce de borde. La misma estructura aparece en la APP levantada,
en la descomposición \(1215=54+9\cdot129\), en la unidad dimensional
estratificada \(1000=729+243+27+1\) y en la lectura de la masa como ruta con
avance visible y enrollamiento contraangular. Esta dualidad enlaza el bloque
\(243\), la bola ternaria perfecta \(V_3(11,2)=243\), el código ternario
puncturado de Golay \([11,6,5]_3\), los canales de
\(\pi,e,\varphi,\alpha\), la bisagra \(744\), la dimensión
\(11=6+5=3+8\) y la involución
\[
(r,Q;R)\longmapsto(Q,r;R^{-1}).
\]
En lenguaje de cuerdas, el valor visible se publica como modo de vibración o
momento, mientras que la memoria de frontera se publica como enrollamiento.
En Mecánica Dimensional estos dos modos aparecen como
\(n_AA_{\mathrm{rad}}\) y \((s_C/6)C^*_{\mathrm{rad}}\): avance visible y
memoria contraangular.

La HMT no toma el continuo como objeto primario. Parte de una célula discreta
con memoria de frontera. La recta real, Moonshine y la física dimensional son
lecturas diferentes de esa misma célula. En consecuencia, la cadena
\[
\text{APP levantada}
\longrightarrow\text{valor y memoria}
\longrightarrow\text{dualidad }T
\longrightarrow243/\text{Golay}
\longrightarrow\text{rutas de masa}
\]
no importa la dualidad desde la teoría de cuerdas: muestra que su estructura
existe en HMT antes de la geometrización continua.

\subsection{Residuo visible y cociente de frontera}

Para todo entero positivo \(n\), defínanse
\[
r_9(n)=1+((n-1)\bmod 9),
\qquad
Q_9(n)=\frac{n-r_9(n)}{9}.
\]
Entonces
\[
n=9Q_9(n)+r_9(n).
\]
La lectura ordinaria conserva el residuo. La lectura enriquecida conserva el
par completo
\[
n\longmapsto(r,Q).
\]
Este par es la estructura mínima de una lectura compacta: una componente se
publica dentro de la rueda y la otra conserva cuántas veces se atravesó su
frontera.

\begin{principle}[Dualidad entre valor visible y memoria de frontera]
La célula HMT admite dos lecturas conjugadas. La primera publica \(r\) como
fase, valor o vibración; la segunda publica \(Q\) como memoria,
enrollamiento o número de cruces de frontera.
\end{principle}

\subsection{Involución de dualidad T}

Sea \(R>0\). El funcional
\[
E_R(r,Q)=\frac{r^2}{R^2}+Q^2R^2
\]
satisface
\[
E_R(r,Q)=E_{1/R}(Q,r),
\]
pues
\[
E_{1/R}(Q,r)
=\frac{Q^2}{(1/R)^2}+r^2(1/R)^2
=Q^2R^2+\frac{r^2}{R^2}.
\]
Por tanto,
\[
(r,Q;R)\longleftrightarrow(Q,r;R^{-1}).
\]
La expresión continua de cuerdas contiene términos de la forma
\[
\frac{n^2}{R^2}+\frac{w^2R^2}{\alpha'^2}.
\]
HMT aporta antes de esa realización continua el residuo visible, la memoria
de frontera y la inversión de escala. La dualidad \(T\) es la proyección
geométrica de esta involución de lectura.

\subsection{APP levantada: modos visibles y modos de memoria}

Para la suma y el producto sobre el mismo soporte se escribe
\[
i+j=9Q^+_{ij}+R^+_{ij},
\qquad
ij=9Q^\times_{ij}+R^\times_{ij}.
\]
La APP levantada es
\[
\mathcal A_9^\sharp=(R^+,Q^+;R^\times,Q^\times),
\]
donde \(R^+,R^\times\) son los modos visibles y
\(Q^+,Q^\times\) los modos de memoria. La discrepancia completa es
\[
ij-(i+j)
=(R^\times_{ij}-R^+_{ij})
+9(Q^\times_{ij}-Q^+_{ij}).
\]
Al sumar sobre \(1\leq i,j\leq9\),
\[
\sum R^+=405,
\quad
\sum R^\times=459,
\quad
\sum Q^+=45,
\quad
\sum Q^\times=174.
\]
De aquí se obtienen
\[
\Delta_{\mathrm{visible}}=459-405=54,
\qquad
\Delta_{\mathrm{memoria}}=174-45=129,
\]
y
\[
\sum ij-\sum(i+j)=2025-810=1215=54+9\cdot129.
\]
La normalización por la rueda produce
\[
\frac{1215}{9}=135=6+129=120+15.
\]
Así, \(54\) es el defecto visible, \(129\) el defecto de memoria,
\(135\) el defecto total normalizado, \(120\) el canal octádico constitutivo
y \(15\) el par interno de la héxada.

\subsection{Completación estratificada de la unidad}

La Extrema y Media Razón Holográfica posee la completación cúbica
\[
10^3=9^3+3\cdot9^2+3\cdot9+1,
\]
es decir,
\[
1000=729+243+27+1.
\]
Los cuatro sumandos representan, respectivamente, el volumen discreto de
dimensión \(3\), la frontera de codimensión \(1\), la frontera de
codimensión \(2\) y el retorno de codimensión \(3\). La corona activa se
refina como
\[
270=243+27=3\cdot9^2+3\cdot9.
\]
La unidad se descompone así en interior, frontera activa estratificada y
retorno. El paso \(9\to10\) abre una capa dimensional completa, y la frontera
de un nivel se convierte en canal de propagación del nivel siguiente.

\subsection{El bloque \texorpdfstring{\(243\)}{243} como frontera y bola ternaria perfecta}

El mismo bloque se obtiene como bola de Hamming ternaria de radio \(2\) y
longitud \(11\):
\[
V_3(11,2)
=\sum_{k=0}^{2}\binom{11}{k}2^k
=1+22+220
=243.
\]
El código ternario perfecto puncturado de Golay tiene parámetros
\[
G_{11}=[11,6,5]_3,
\qquad
|G_{11}|=3^6=729,
\]
y satisface
\[
729\cdot243=3^{11}.
\]
Por tanto, \(G_{11}\) tesela \(\mathbb F_3^{11}\) mediante bolas ternarias de
radio \(2\). El número \(243\) es simultáneamente
\[
3\cdot9^2,
\qquad
V_3(11,2),
\qquad
3^5,
\qquad
1+22+220.
\]
Esta identidad une frontera dimensional, corrección perfecta, dimensión
\(11\), incidencia de Witt--Golay y la lectura estructural de teoría M.

\subsection{Descomposición interna del bloque \texorpdfstring{\(243\)}{243}}

La bola admite la descomposición
\[
243=1+22+90+120+10.
\]
Aquí \(1\) es el centro y retorno; \(22=2\cdot11\), el primer cascarón de
Hamming y la lectura electrónica del bloque; \(90=6\binom62\), el canal hexádico;
\(120=8\binom62\), el canal octádico; y \(10=9+1\), el cierre de la
monodromía \(9\to1\). Así, retorno, componente \(22\), canales \(90/120\) y
monodromía quedan contenidos en una misma bola perfecta.

\subsection{Canales de constantes derivados del bloque \texorpdfstring{\(243\)}{243}}

Con \(B=243\), se obtienen los canales
\[
C_\alpha=3B=729,
\qquad
C_e=B+27+1=271,
\]
\[
C_\pi=B+72=315,
\qquad
C_\varphi=B-81-1=161.
\]
El canal de \(\alpha\) es el bloque triple \(243\); el de \(e\), el bloque
\(243\) con la frontera \(27\) y el retorno; el de \(\pi\), el bloque con el
canal \(72\); y el de \(\varphi\), el bloque después de retirar el
subcuadrado \(81\) y el retorno.
La bisagra se reconstruye como
\[
C_\pi+C_e+C_\varphi-3
=315+271+161-3
=744
=729+15.
\]
La traza reducida de los tres caracteres combina, por tanto, el bloque triple
\(243\) con el par interno de la héxada.

\subsection{Dimensión \texorpdfstring{\(11\)}{11} y lecturas compatibles}

El código ternario perfecto puncturado produce
\[
11=6+5,
\]
donde \(6\) es la dimensión de \(G_{11}\) y \(5\) la profundidad de
corrección \(3^5=243\). La misma dimensión se descompone como
\[
11=3+8,
\]
donde \(3\) corresponde a TRIT y \(8\) a la octada interna. En conjunto,
\[
11=6+5=3+8,
\qquad
12=11+1.
\]
La lectura de dimensión \(11\) coordina así código y corrección, TRIT y
octada, y la identidad \(729\cdot243=3^{11}\).

\subsection{Realización dimensional de la dualidad}

La acción local de una ruta se expresa como
\[
E_0(P)=n_A(P)A_{\mathrm{rad}}
+\frac{s_C(P)}{6}C^*_{\mathrm{rad}}.
\]
El primer término es el modo de avance o vibración visible; el segundo, el
modo de frontera o enrollamiento contraangular. La masa local es
\[
m_{\mathrm{HMT}}(P)=m_e\exp(E_0(P)).
\]
La torre refinada añade la memoria de frontera:
\[
E(P)=E_0(P)+k(P)\Delta_4
+\nu_{120}(P)s_{120}
+\nu_{270}(P)\zeta_{270},
\]
donde \(\Delta_4\) es la torsión mínima, \(s_{120}\) el selector
hexada--octada y
\[
\zeta_{270}=135\Delta_4^2
\]
la memoria cuadrática de la corona. La dualidad abstracta de ruta es
\[
E_{A,C}(n,w)=nA_{\mathrm{rad}}+wC^*_{\mathrm{rad}},
\qquad
E_{A,C}(n,w)=E_{C,A}(w,n).
\]
En el sector físico \(A\) y \(C^*\) están polarizados, y el registro genealógico selecciona
la realización. La masa es la exponencial de una ruta de cierre con avance
visible, enrollamiento de frontera y torsión.

\subsection{Compatibilidad entre hojas y cierre del registro genealógico}

En una cuerda cerrada, la compatibilidad entre los modos izquierdo y derecho
se expresa esquemáticamente por
\[
N_L-N_R=nw.
\]
En la lectura HMT, el miembro izquierdo mide la diferencia de excitación
entre hojas, mientras el derecho mide el acoplamiento entre la fase visible y
la memoria de frontera. El ajuste de niveles constituye así un cierre del
registro genealógico bidireccional. En las lecturas \(CP\), \(T\) y \(CPT\), una asimetría
de orientación se compensa mediante la inversión bidireccional para cerrar el
registro.

\subsection{Dualidad T, modularidad y función J}

En el toro complejo,
\[
q=e^{2\pi i\tau},
\]
y la transformación modular
\[
\tau\longmapsto-\frac1\tau
\]
intercambia sus ciclos. La dualidad \(T\) intercambia radios grandes y
pequeños; ambas realizaciones publican el intercambio entre el ciclo visible
y el ciclo de frontera. La función modular \(J\), la acción de
Sommerfeld y la dualidad \(T\) constituyen cierres distintos de una misma
estructura primaria de frontera con memoria.

\subsection{Consecuencias dimensionales para fotón, electrón, masa y estadística}

El fotón observable se asocia al modo central del TPK:
\[
L_{\mathrm{TPK}}=3(I-P_0),
\qquad
\operatorname{Spec}(L_{\mathrm{TPK}})=\{0,3,3\}.
\]
El modo cero es una propagación nula extendida sin memoria compactificada
propia. El electrón aparece como cierre compacto de los modos nulos
conjugados,
\[
e^-=\text{cierre compacto de }P_++P_-,
\]
y constituye el enrollamiento físico mínimo con memoria estable. La masa
mide el coste de estabilización de esa memoria compacta. La exclusión de
Pauli se interpreta como saturación del registro genealógico: dos fermiones iguales no
ocupan la misma frontera porque su cierre de enrollamiento ya está saturado.
Los bosones publican modos compartibles de memoria desplegada, mientras los
fermiones publican memoria compactificada.

\subsection{Identidades verificables de la construcción}

La construcción satisface conjuntamente las identidades
\begin{enumerate}
\item \(1000=729+243+27+1\);
\item \(V_3(11,2)=1+22+220=243\);
\item \(729\cdot243=3^{11}\);
\item \(243=1+22+90+120+10\);
\item \(C_\alpha,C_e,C_\pi,C_\varphi=729,271,315,161\);
\item \(315+271+161-3=744\);
\item \(1215=54+9\cdot129\);
\item \(E_R(r,Q)=E_{1/R}(Q,r)\);
\item \(11=6+5=3+8\), \(12=11+1\), \(26=2+24\) y \(10=2+8\).
\end{enumerate}

La dualidad \(T\) se reconoce, por tanto, como la realización geométrica del
intercambio entre residuo y cociente. El enrollamiento no se introduce
primero como propiedad externa de una cuerda compacta, sino como memoria de
frontera de la descomposición \(n=9Q+r\). El bloque \(243\) reúne esta
dualidad con Golay ternario, dimensión \(11\), canales \(90/120\), generación
de constantes, rutas de masa y teoría M. En su forma más general, HMT lee la
masa, la cuerda y la dimensión como publicaciones de una misma frontera con
memoria.
